\subsection{6. Feldtyp: was tatsächlich erzwungen
ist}\label{feldtyp-was-tatsuxe4chlich-erzwungen-ist}

Für ein \textbf{lokales, ableitungsfreies Bilinear zweier gleichhändiger
Weylfelder}, das alle 64 inneren Quellenmarken unverändert als
unabhängige innere Komponenten und genau den gegebenen antisymmetrischen
Tensor W verwendet, gilt

\codeword{(1/2,0) ⊗ (1/2,0) = (1,0) ⊕\allowbreak{} (0,0)}.

Die skalare Epsilon-Kontraktion zusammen mit antisymmetrischem W ist im
gemeinsamen Index symmetrisch und verschwindet wegen
Grassmann-Antikommutation. Der symmetrische Spinortensor, also der
(1,0)-Kanal, bleibt nichtverschwindend. Im genannten Vertrag ist dies
korrekt. Ein zusätzliches unabhängiges gleichgeladenes Hilfsdublett
erlaubt stattdessen wieder die skalare Kontraktion.

Die Zahlen 180/128 beziehungsweise 60/256 zählen \textbf{komplexe lokale
Feldkomponenten im jeweiligen direkten Tensorprodukt-Ansatz}. Sie zählen
nicht automatisch zusätzliche unabhängige physische Oszillatoren oder
propagierende Freiheitsgrade. Diese hängen von Kinetik,
Nebenbedingungen, positiver Energie, Teilchen-/Antiteilchenstruktur und
dem tatsächlichen Feldadapter ab. Deshalb folgt aus der Multiplikation
mit zwei oder drei kein allgemeiner Ausschluss jeder relativistischen
Lesart mit den nativen Modenzahlen.

Vor allem folgt daraus keine vorgeschriebene Reihenfolge »erst räumliche
Skalierung, dann Feldwörterbuch«. Lokale Darstellungstypen und mögliche
Kinetik können vor einer Skalierungsrechnung geprüft werden. Welche
Feldkomponenten tatsächlich aus dem räumlichen Quellprozess entstehen,
muss anschließend gemeinsam mit dem Adapter untersucht werden.

\subsubsection{Konkreter Gegencheck jenseits der ableitungsfreien
Klasse}\label{konkreter-gegencheck-jenseits-der-ableitungsfreien-klasse}

Der allgemeine Satz »ein Vektor verlangt ψ†ψ und scheitert deshalb an
der Ladung« ist falsch, sobald Ableitungen zugelassen werden.
Beispielsweise ist

\codeword{J\textasciicircum{}A\_mu = Σ\_IJab M\textasciicircum{}A\_IJ ε\_ab ψ\_Ia ↔∂\_mu ψ\_Jb}

ein Lorentzvektor mit Ladung −2. Die Ableitungsantisymmetrisierung macht
ihn bei antisymmetrischem M nichtverschwindend. Für M=ε und zwei inneren
Marken liefert die exakte Grassmann-Jetrechnung vier nichtverschwindende
Monome mit Koeffizienten +2,−2,−2,+2. Ein entsprechend geladener
Vektorvermittler könnte algebraisch durch
\codeword{b†\_mu J\textasciicircum{}mu + h.c.} koppeln, ohne
Ladungsverletzung.

Das ist \textbf{nur ein Gegenbeispiel gegen den überbreiten Ausschluss},
keine native Lösung: Es fügt eine Ableitung und einen Feldtyp hinzu; bei
kanonischen 4D-Felddimensionen ist es ein höherdimensionaler
Kopplungsterm. Gesunde Kinetik, Quellenherkunft, Skalierung und T1-T8
werden dadurch nicht bewiesen.

\subsection{7. Konsequenz}\label{konsequenz}

Die neuen Quellen verkleinern und präzisieren den endlichen
N=3-Operationsraum. Sie schließen weder die physische Verfügbarkeit der
Operationen noch den Lorentzadapter. Die sachlich tragfähige Integration
übernimmt die Zerlegung, Casimire, 14/16/7 im jeweils erklärten Vertrag
sowie den eingeschränkten Feldtypsatz; sie ersetzt die überbreiten
Ausschlüsse und die behauptete Programmumkehr durch explizite
Bedingungen.

\begin{center}\rule{0.5\linewidth}{0.5pt}\end{center}

\section{Anhang B: unabhängige Krylov- und
Feldprüfung}\label{anhang-b-unabhuxe4ngige-krylov--und-feldpruxfcfung}

\section{Unabhängige Prüfung: Krylov-Verzweigung und
Feldwörterbuch}\label{unabhuxe4ngige-pruxfcfung-krylov-verzweigung-und-feldwuxf6rterbuch}

Arbeitsstand 2026-09-15. Geprüft wird die Anlage
\codeword{7380c383-f364-4ef8-8066-24aa4d180d85} gegen den Vertrag
\codeword{universalraum-native-operations-ground-response-20260915}.
Fremde Quellen und Ausgaben wurden nicht verändert. Die hier angegebenen
Matrix- und Graphbefunde beziehen sich auf denselben gepinnten Tensor W
mit SHA-256
\codeword{3f00a0892157a691e4ef26920c3702772c7bf76a250fad98ff04a079bb64b763}.

\subsection{1. Der neue Nichtschluss ist exakt
bestätigt}\label{der-neue-nichtschluss-ist-exakt-bestuxe4tigt}

Für \codeword{v\_n=(T†)\textasciicircum{}n F} habe ich
\codeword{u=T T† v\_2=T v\_3} unmittelbar aus den 480 nativen Paaren neu
kontrahiert. Das materialisiert nicht die 15.252.960 Komponenten von v₃.
Es arbeitet mit 108.240 Komponenten von v₂ und 293.280
nichtverschwindenden Komponenten des Ergebnisses u. Es wurden 45.771.360
zulässige Erzeugungsübergänge und 139.289.760 anschließende
Entnahmebeiträge exakt mit ganzzahliger Fermion-Parität und
Boson-Multiplizität zusammengeführt.

Die unabhängige Rechnung liefert

\[\adjustbox{max width=\linewidth}{$\displaystyle 
\|v_2\|^2=439680,\quad
\langle v_2,u\rangle=575078400,\quad
\|u\|^2=752194252800.
$}\]

Damit folgen rein rational

\[\adjustbox{max width=\linewidth}{$\displaystyle 
\alpha=\frac{299520}{229},\qquad
w_2=u-\alpha v_2,\qquad
\|w_2\|^2=\frac{5001523200}{229}>0,
$}\]

und

\[\adjustbox{max width=\linewidth}{$\displaystyle 
\frac{\|w_2\|^2}{\nu_3}=\frac{1809}{47632}.
$}\]

Der Ein-Vektor-pro-Bosonlage-Ansatz ist somit nicht invariant. Da T die
native Symmetrie erhält und v₂ und u Singuletts sind, liefert w₂ einen
zweiten unabhängigen Singulettvektor in derselben Bosonlage. Das ist ein
Strukturresultat, kein physikalischer Mehrteilchen- oder
Raumzeitnachweis.

\subsection{2. Das mitgelieferte Ground-JSON ist veraltet und teilweise
falsch}\label{das-mitgelieferte-ground-json-ist-veraltet-und-teilweise-falsch}

Die aktuelle Pythonquelle hat SHA-256
\codeword{8ae0f2f103ba0cb0a5cb832681a33130b6756fe627395368b6be635587a2eb02}.
Das vorhandene \codeword{native\_ground\_state.json} nennt dagegen
\codeword{046ffa3f41d9dd365e2c9a199e938734bf735d762f23373a3f85b541c3da30fe}
und enthält noch

\[\adjustbox{max width=\linewidth}{$\displaystyle 
\|w_2\|^2_{\rm alt}=1145348812800
=229^2\|w_2\|^2.
$}\]

Sein \codeword{closure\_defect=1.5226768996658837} ist schon als
relative orthogonale Norm unmöglich. Das daraus abgeleitete
Ritz-Ergebnis −1,18826198 ist ungültig. Der spätere Anlagentext
korrigiert den Normfehler, aber die JSON-Belegkette wurde noch nicht
entsprechend neu erzeugt. Das grüne \codeword{status: PASS} dieser
historischen Ausgabe darf deshalb nicht mit dem korrigierten Text
zusammengezogen werden.

Die korrigierte Norm und die folgenden kleinen Matrizen wurden hier
unabhängig reproduziert. Die ganze fremde große Ausführung wurde
ausdrücklich nicht erneut gestartet.

\subsection{3. Die sechste Dezimalstelle ist richtig - aber nur für die
eine zusätzliche
Richtung}\label{die-sechste-dezimalstelle-ist-richtig---aber-nur-fuxfcr-die-eine-zusuxe4tzliche-richtung}

Die neue numerische Diagonalisierung ergibt:

{\def\LTcaptype{none} % do not increment counter
\begin{longtable}[]{@{}
  >{\raggedright\arraybackslash}p{(\linewidth - 6\tabcolsep) * \real{0.2000}}
  >{\raggedleft\arraybackslash}p{(\linewidth - 6\tabcolsep) * \real{0.2667}}
  >{\raggedleft\arraybackslash}p{(\linewidth - 6\tabcolsep) * \real{0.2667}}
  >{\raggedleft\arraybackslash}p{(\linewidth - 6\tabcolsep) * \real{0.2667}}@{}}
\toprule\noalign{}
\begin{minipage}[b]{\linewidth}\raggedright
g/Δ
\end{minipage} & \begin{minipage}[b]{\linewidth}\raggedleft
vierdimensionale Ritz-Matrix
\end{minipage} & \begin{minipage}[b]{\linewidth}\raggedleft
plus w₂
\end{minipage} & \begin{minipage}[b]{\linewidth}\raggedleft
Änderung
\end{minipage} \\
\midrule\noalign{}
\endhead
\bottomrule\noalign{}
\endlastfoot
1/20 & −1,0942308394409612 & −1,0942318710142458 & −1,0315733·10⁻⁶ \\
1/40 & −0,29535718178032033 & −0,29535720496683315 & −2,3186513·10⁻⁸ \\
1/100 & −0,04789424798601009 & −0,04789424801311775 &
−2,7107663·10⁻¹¹ \\
\end{longtable}
}

Das sind numerische Eigenwerte exakter kleiner Ritz-Matrizen, keine nach
außen gerundeten vollständigen Spektraleinschließungen. Durch
Eliminieren der einzigen neuen w₂-Koordinate bekommt die v₃-Diagonale
die exakt bestimmte energieabhängige Korrektur

\[\adjustbox{max width=\linewidth}{$\displaystyle 
-\frac{g^2(1809/47632)}{2\Delta-E}.
$}\]

Das erklärt die kleine Änderung durch diesen einen Zweig. Es sagt nichts
über die ausgelassenen höheren Bosonlagen.

\subsubsection{Eigene einfache Folgeherleitung: Das vollständige
Ritz-Residuum}\label{eigene-einfache-folgeherleitung-das-vollstuxe4ndige-ritz-residuum}

Sei ψ₃ der normierte Ritz-Eigenvektor in \codeword{span(v₀,…,v₃)}, und
c₃ seine Komponente entlang \codeword{v₃/√ν₃}. Außerhalb dieses Raums
führt Hψ₃ genau in zwei orthogonale Richtungen: v₄ und w₂. Deshalb gilt
exakt

\[\adjustbox{max width=\linewidth}{$\displaystyle 
\|(H-E_{\rm Ritz})\psi_3\|^2
=g^2|c_3|^2\frac{\nu_4+\|w_2\|^2}{\nu_3}.
$}\]

Hier wird \codeword{ν₄=952296652800} aus dem früher vollständig
geprüften Momentbeleg übernommen; die vierte Ordnung wurde in dieser
Runde nicht erneut enumeriert.

Bei g/Δ=1/20 folgt numerisch ein Residuum von rund \textbf{0,373063 Δ}.
Der w₂-Zweig trägt weniger als \textbf{23 Millionstel} des quadrierten
Residuums. Die große ausgelassene Richtung ist v₄, nicht w₂. Aus dem
winzigen Nichtschlussanteil darf daher keine praktisch vollständige
Konvergenz abgeleitet werden.

Außerdem liegt die neue Fünf-Zustands-Ritzenergie mindestens
\textbf{0,0354 Δ über} der bereits bewiesenen Obergrenze des wahren
Grundzustands. Das ist eine nachweislich noch relevante Lücke,
unabhängig von einer Konvergenzschätzung.

\subsection{4. Der bereits bewiesene native Grundsatz bleibt
gültig}\label{der-bereits-bewiesene-native-grundsatz-bleibt-guxfcltig}

Der ältere, gepinnte Beweis lautet: Für Δ\textgreater0 und
0\textless\textbar g\textbar/Δ≤1/20 besitzt der unveränderte
μ=0-Hamiltonoperator einen eindeutigen globalen Grundzustand in N=64; er
ist ein Spin(10)×SU(4)-Singulett. Bei 1/20 bestehen außerdem positive
Lücke und die bekannten Energie-, Besetzungs- und Polschranken.

Die neue Ausführung benutzt schwächere Sektor-Untergrenzen und weniger
Ritz-Vektoren. Dass diese schwächere Methode bei 1/20 nicht genügt, ist
weder eine Widerlegung noch ein Wiederöffnen des früheren Satzes.
Insbesondere sind die vorgeschlagenen Ziele \codeword{E<−1,2Δ} bzw.
\codeword{E<−1,1625Δ} unmöglich, weil bereits \codeword{E₀>−1,158089Δ}
bewiesen ist. Bessere Schranken für die Konkurrenzsektoren sind
erforderlich, nicht eine unrealistisch tiefere Energie.

Die neuen Zahlen \codeword{Nb≈0,84211266}, \codeword{Zh≈0,97368398} und
\codeword{εh≈0,03107309Δ} gehören zum nichtstationären K₃-Ritz-Zustand.
Nb liegt sogar knapp unter der früher bewiesenen Schranke
\codeword{Nb(Ω)>0,842846}. Zh bezeichnet die gesamte Entnahmenorm dieser
Näherung, nicht das Gewicht der isolierten niedrigen Linie. Ein erstes
Moment auf einem Ritz-Zustand ist keine neu bestimmte Polenergie auf Ω.

Offen bleiben der vollständige Ω-Vektor, die physische Auswahl des
μ=0-Vertrags und die Herkunft ausführbarer räumlicher Operationen.

\subsection{5. Feldwörterbuch: bestätigte Algebra und notwendige
Einschränkungen}\label{feldwuxf6rterbuch-bestuxe4tigte-algebra-und-notwendige-einschruxe4nkungen}

Frisch exakt geprüft wurden:

\begin{itemize}
\tightlist
\item
  Der Trägergraph ist zusammenhängend, 15-regulär und dreiecksfrei.
\item
  Der Fünf-Zyklus \codeword{16→45→0→30→37→16} liegt wirklich im
  Trägergraphen.
\item
  Für jeden der 60 nativen Kanäle verschwindet der gleichhändige
  Ein-Kopie-Weyl-Skalaransatz \codeword{M\_A⊗ε} als Grassmann-Bilinear.
\item
  Bei Dirac-Kernen sind C, Cγ⁵ und alle vier Cγ\^{}μγ⁵ antisymmetrisch;
  alle vier Cγ\^{}μ und alle sechs Cσ\^{}\{μν\} sind symmetrisch.
\item
  Die 43 getesteten nichttrivialen Gradierungen ergeben genau die
  angegebenen Kanalzahlen 24/36, 40/20 bzw. 32/28. Die
  Stabilisatordimensionen 36, 52 und 28 wurden diesmal durch exakt
  diagonale reelle Gram-Matrizen der Kommutatorbedingungen bestätigt,
  nicht durch einen SVD-Schwellwert.
\end{itemize}

Der Fünf-Zyklus allein schließt nur eine diagonale
±-Händigkeitszuweisung aus, die an jedem Paar entgegengesetzte
Händigkeiten verlangt. Er allein ist kein basisunabhängiger Ausschluss
und kein Ausschluss zusätzlicher Feldkomponenten.

\subsubsection{Eigener stärkerer Satz: auch ein Basiswechsel repariert
den Ein-Kopie-Vektoransatz
nicht}\label{eigener-stuxe4rkerer-satz-auch-ein-basiswechsel-repariert-den-ein-kopie-vektoransatz-nicht}

Sei Γ eine hermitesche Händigkeit auf dem ursprünglichen
64-dimensionalen internen Raum. Für einen ausschließlich
entgegengesetzt-händigen, unveränderten Paartensor müsste

\[\adjustbox{max width=\linewidth}{$\displaystyle 
\Gamma^T M_A+M_A\Gamma=0\quad\text{für alle }A
$}\]

gelten. Adjungieren und Einsetzen zeigt

\[\adjustbox{max width=\linewidth}{$\displaystyle 
[\Gamma,M_A^\dagger M_B]=0\quad\text{für alle }A,B.
$}\]

Die native Faktorisierung wurde exakt rekonstruiert:

\[\adjustbox{max width=\linewidth}{$\displaystyle 
M_{k,a}=S_k\otimes C_a,\quad
\sum_k S_k^\dagger S_k=5I_{16},\quad
\sum_a C_a^\dagger C_a=3I_4.
$}\]

Daher muss Γ sowohl mit allen \codeword{S\_k†S\_l⊗I₄} als auch mit allen
\codeword{I₁₆⊗C\_a†C\_b} kommutieren. Die von den ersten Produktmatrizen
erzeugte Algebra ist die volle M₁₆; die zweite ist M₄. Hierfür enthält
die Prüfung vollständige Rangzertifikate von ganzzahligen Matrixwörtern
modulo 101: \textbf{256 unabhängige Spinorwörter bis Länge zwei und 16
unabhängige Farbprodukte bis Länge eins}. Voller Rang modulo einer
Primzahl beweist vollen Rang über den komplexen Zahlen; es wird
ausdrücklich keine Rangdefizienz von einem endlichen Körper übertragen.

Folglich ist Γ skalar. Die ursprüngliche Gleichung erzwingt dann Γ=0.
Insbesondere gibt es \textbf{keine hermitesche Involution Γ²=I} mit der
verlangten Eigenschaft. Diese Schlusskette schließt den
Ein-Weyl-pro-Label-Vektoradapter jetzt unabhängig von einer Wahl der
internen Basis aus.

Der Satz betrifft unveränderte W-Kanäle und unveränderten
64-dimensionalen Träger. Er schließt \textbf{nicht} Diracfelder pro
Label, zusätzliche Kopien, Ableitungsadapter oder andere ausdrücklich
erweiterte Modelle aus.

\subsection{6. Die minimale Erweiterung darf nicht aus Versehen mit
ausgeschlossen
werden}\label{die-minimale-erweiterung-darf-nicht-aus-versehen-mit-ausgeschlossen-werden}

Zwei explizite Gegenkontrollen wurden auf allen 60 Kanälen neu
gerechnet:

\begin{enumerate}
\def\labelenumi{\arabic{enumi}.}
\tightlist
\item
  \textbf{Vierkomponenten-Diracfelder pro Label:} \codeword{M\_A⊗Cγ⁰}
  ist nichtverschwindend und antisymmetrisch im gesamten
  Grassmann-Index. Es enthält zusätzliche linke und rechte
  Freiheitsgrade pro ursprünglichem Label. Deshalb widerspricht es dem
  Ein-Kopie-Satz nicht. Eine chirale Standardmodell-Entstehung folgt
  daraus nicht.
\item
  \textbf{Unabhängige Zweier-Erweiterung:}
  \codeword{M\_A⊗ε\_aux⊗ε\_Lorentz} ist nichtverschwindend und
  antisymmetrisch. Der skalare Kanal wird möglich, weil die zusätzlich
  antisymmetrische interne Zweierform die Symmetrie des internen
  Paarfaktors umkehrt. Das ist die bereits bekannte minimale Reparatur
  im deklarierten Tensorprodukt-Ansatz, keine aus der Quelle
  hergeleitete Verdopplung.
\end{enumerate}

Eine solche Zweier-Erweiterung macht auch den Trägergraphen durch
\codeword{M\_A⊗σ\_x} bipartit: \codeword{Γ=I₆₄⊗diag(1,−1)}
antikommutiert exakt mit allen so erweiterten Matrizen. Damit ist
sichtbar, welche zusätzliche Ressource den ursprünglichen Ausschluss
überwindet.

``Der Vermittler kann nie Skalar sein'' muss folglich auf den
unveränderten Ein-Kopie-Ansatz begrenzt werden. Ebenso fixiert Schur nur
die interne Matrixstruktur innerhalb eines irreduziblen Blocks; es
beweist nicht die Einzigartigkeit aller möglichen Lorentz- und
Ableitungsterme. Eine gesunde Kinetik, Nebenbedingungen, Eichherkunft
und ein dynamischer Spin-2-Sektor sind hier nicht konstruiert.

